% conclusion.tex — Full proposal draft, Week 8

\section{Conclusion}

This project tests whether a brief, portable brain recording session, calibrated on naturalistic movie-watching, can serve as a window into whole-brain neural dynamics during structured cognitive tasks and into individual differences in attentional control. The preliminary results from the movie-watching training phase show that the aPCR model recovers large-scale brain organization from fNIRS surface signals alone, with significant prediction in 56 of 122 brain regions and strong preservation of inter-network connectivity structure ($r = 0.617$). The dorsal attention network (DAN) shows the strongest proportional performance (11 of 14 ROIs, 79\%), with the frontoparietal control network and default mode network also significantly predicted, providing a foundation for the cross-task individual differences analyses. If those analyses show that predicted brain patterns correlate with reaction time variability and task performance, the result would be a candidate method for conducting whole-brain-equivalent neurocognitive assessment without a scanner, with direct implications for making attention research accessible to the populations and settings where it is needed most.
